\documentclass[12pt]{article}
\usepackage[T1]{fontenc}
\usepackage[utf8]{inputenc}
\usepackage{lmodern}
\usepackage{textcomp}
\usepackage{enumitem}
\usepackage{geometry}
\geometry{margin=1in}
\begin{document}
\begin{center}
Answer Key - Business Administration - Graduate Level Assessment 24 \
\small (Answers are ordered exactly as in the question paper.)
\end{center}

\section*{Multiple Choice}
\begin{enumerate}[label=\arabic*.]
\item \textbf{Answer:} A
\\ \textit{Options:} A) 28\%; B) 15\%; C) 22\%; D) 10\%; E) 30\%
\\ \textit{Note:} The correct answer is 28\% because if a sofa is marked up by 20\% based on its selling price, the cost represents approximately 83.33\% of the selling price, leading to a markup of about 28\% when calculated based on the cost.
\item \textbf{Answer:} A
\\ \textit{Options:} A) \$840.00; B) \$729.60; C) \$512.00; D) \$691.20; E) \$768.00
\\ \textit{Note:} The correct answer of $840.00 is derived from applying a 20% declining balance depreciation method over three years, where the book value is calculated as follows: Year 1 depreciation is $240 (20\% of $1200), leaving a book value of $960; Year 2 depreciation is $192 (20% of $960), leaving a book value of $768; and Year 3 depreciation is $153.60 (20\% of $768), resulting in a final book value of $840.00.
\item \textbf{Answer:} B
\\ \textit{Options:} A) \$200.00; B) \$176.69; C) \$165.50; D) \$150.75; E) \$135.20
\\ \textit{Note:} The correct answer of $176.69 is derived from applying the formula for simple interest, which is calculated as \(I = P \times r \times t\), where \(P = 1285\), \(r = 0.055\), and \(t = 2.5\) years, leading to an interest amount of approximately $176.69.
\item \textbf{Answer:} A
\\ \textit{Options:} A) \$14.95; B) \$15.95; C) \$16.00; D) \$14.50; E) \$15.05
\\ \textit{Note:} The correct answer is $14.95 because it represents the previous year's monthly charge, calculated by taking the current average bill of $17.05 and dividing it by 1.10 to account for the 10\% increase.
\item \textbf{Answer:} A
\\ \textit{Options:} A) \$1.13; B) \$3.59; C) \$2.70; D) \$2.00; E) \$2.40
\\ \textit{Note:} The correct answer is $1.13 because the cost can be calculated by dividing the selling price by 1 plus the markup percentage (1 + 0.35), resulting in $3.24 / 1.35 = \$1.13.
\end{enumerate}

\section*{True / False}
\begin{enumerate}[label=\arabic*.]
\setcounter{enumi}{5}
\item \textbf{Answer:} True
\\ \textit{Options:} A) True; B) False
\\ \textit{Note:} The job characteristics model, developed by Hackman and Oldham, emphasizes core job dimensions such as skill variety, task identity, task significance, autonomy, and feedback as crucial elements for enhancing employee motivation, while not explicitly including 'reward' as a job design characteristic.
\item \textbf{Answer:} True
\\ \textit{Options:} A) True; B) False
\\ \textit{Note:} The statement is true because the total cost of producing one item is the sum of the fixed cost of $4 and the marginal cost of $3.45, which together equal \$7.45.
\item \textbf{Answer:} False
\\ \textit{Options:} A) True; B) False
\\ \textit{Note:} Transformative working practices prioritize flexibility, collaboration, and adaptability over rigid structures and fixed roles to foster innovation and improve productivity.
\item \textbf{Answer:} False
\\ \textit{Options:} A) True; B) False
\\ \textit{Note:} The statement is false because the number of competitors in the market directly influences the level of service surrounding a product, as increased competition typically drives businesses to enhance their service offerings to attract and retain customers.
\item \textbf{Answer:} False
\\ \textit{Options:} A) True; B) False
\\ \textit{Note:} Belbin (1981) actually characterizes the Monitor Evaluator role as someone who is analytical, objective, and tends to be more introverted, rather than extroverted and enthusiastic.
\end{enumerate}

\section*{Long Form Response}
\begin{enumerate}[label=\arabic*.]
\setcounter{enumi}{10}
\item \textbf{Answer:} The selling price of the piano, set at $22,950, represents 127.5% of its cost value of $18,000. This percentage indicates a markup that not only covers the cost of the piano but also contributes to profit margins essential for business sustainability. The implication of pricing at this level suggests a strategic approach where businesses can leverage perceived value and quality to justify higher prices, potentially appealing to a market segment willing to pay a premium. Understanding this relationship aids businesses in developing pricing strategies that align with market conditions and consumer expectations while effectively managing costs to maximize profitability.
\\ \textit{Note:} correct because it effectively calculates the selling price as a percentage of the cost, revealing the markup necessary for profit while demonstrating how strategic pricing can attract customers who perceive higher value, thereby ensuring business sustainability.
\item \textbf{Answer:} Research methods that utilize predetermined, standardized questions often employ non- probability sampling techniques, which can significantly influence the validity and generalizability of findings in business administration. These methods, such as surveys and structured interviews, allow researchers to gather quantitative data efficiently, facilitating comparisons across diverse participant responses. However, the reliance on non-probability sampling can introduce biases, as it may not accurately represent the broader population, limiting the applicability of results. Furthermore, while standardized questions enhance reliability by ensuring consistency, they may also restrict the depth of responses, potentially overlooking nuanced insights that could inform strategic business decisions. Ultimately, understanding these characteristics and their implications is crucial for effectively interpreting research outcomes in the business context.
\\ \textit{Note:} correct because it identifies that predetermined, standardized questions enhance data collection efficiency and comparability in research, but highlights the potential biases introduced by non-probability sampling, which can impact the validity and generalizability of the findings in business administration.
\end{enumerate}

\vfill
\noindent\textit{End of Answer Key}
\end{document}